\noindent
\subsection{WP4: XR Platform Development and UX Design}

\textbf{Leading Partner:} Arts et Métiers (CAPLAB / Laval Virtual)

\textbf{Objectives and Concept:}
This Work Package integrates advanced XR (VR/AR) and 3D technologies to create a new use-case for cultural heritage and tourism. By utilizing CAPLAB's 4D reality capture capabilities (motion, gesture, sound, physiological data), the project aims to preserve "artisanal intelligence" alongside physical structures. Moving the technology from industrial and health sectors into heritage increases its societal and market relevance, and allows for technological validation (increasing TRL) in complex, outdoor environments.

\textbf{Deliverables / Schedule (from Gantt chart):}
\begin{itemize}
    \item \textbf{D4.1 Development of a VR-based "Virtual Tour" engine (Unreal Engine/Unity):} Year 1 (Month 10-12) \& Year 2 (Month 1-3). This task involves building the core 3D rendering architecture capable of handling high-fidelity photogrammetry datasets smoothly. The engine will focus on optimizing lighting, texture loading, and spatial audio to maximize immersion without requiring industrial-grade hardware from the end users.
    \item \textbf{D4.2 User Experience (UX) research and prototype testing with focus groups:} Year 2 (Month 1-3). The team will conduct iterative usability testing with diverse demographic groups to identify pain points in virtual navigation. The goal is to ensure the interface is highly intuitive, accessible, and inclusive for both digital natives and older audiences.
    \item \textbf{D4.3 Development of an AR-based "On-Site Time Machine" module (ruin reconstruction via mobile application):} Year 2 (Month 4-6). This task focuses on augmenting the physical visitor experience by overlaying digital reconstructions onto real-world ruins in real-time. It requires precise spatial tracking and geolocation integration to seamlessly blend the 3D models with the user's camera feed.
\end{itemize}